% !PRINTED 92
\weberpage{92}
%<<CONT

condition contains a new one of the quantities \(y\), are independent of one another. \gesperrt{If one of these conditions is not fulfilled, the given system of equations has no solution.}

If in the system of equations (1) we regard the \(x\)'s not as unknowns but as variables, then the \(y\)'s too will be variable quantities. Under this interpretation we call the system (1) a \gesperrt{linear substitution}, insofar as it mediates the passage from one system of variables to another. Only when the determinant \(A\), which we now call the \gesperrt{determinant of the substitution}, is different from zero can the \(y\)'s be regarded as independent variables. If this is the case, the system (5) gives the representation of the variables \(x\) in terms of the \(y\)'s, or the \gesperrt{substitution inverse} to (1).


\begin{center}
{\large \S.\ 27.}

\vspace{0.5em}
{\large\gesperrt{Multiplication of Determinants.}}
\end{center}

The theorem which we still have to prove teaches how one can represent the product of two determinants with the same number of rows by a single determinant with the same number of rows. One is led to it most simply by considering the solution of two systems of linear equations.

Let now the coefficients $a_i^{(k)}, b_i^{(k)}$, when $i,k$ run through the sequence of numbers $1,2,\ldots,n$, be arbitrary \gesperrt{variable quantities}; likewise let the quantities $x_i,y_i,z_i$ be bound only by the relations
\begin{align}
 \sum_i a_i^{(k)}x_i&=y_k,\tag{1}\\
 \sum_k b_h^{(k)}y_k&=z_h.\tag{2}
\end{align}
If the problem is now posed of determining the variables $x$ by the variables $z$, this problem can be solved in two ways. By § 26, (4) one may solve equations (1) with respect to $x$, equations (2) with respect to $y$, and insert the latter expressions into the former. If $A,B$ denote the determinants of the two systems of equations, that is,
%CONT>>
